\documentclass[a4paper]{article}
%\usepackage{exercise}
%um nur aufgaben zu zeigen
\usepackage[noanswer]{exercise} 
\usepackage{../images/preamble}
\usepackage{rotating}
\usetikzlibrary{decorations.pathmorphing}
\usetikzlibrary{decorations.markings}
\usetikzlibrary{arrows}
\usetikzlibrary{shapes.geometric}
\newcommand{\midarrow}{\tikz \draw[-triangle 90] (0,0) -- +(.02,0);}
\usepackage{xcolor}
%\usepackage{draftwatermark}
%\SetWatermarkText{\textsc{Draft 2}}
%\SetWatermarkScale{3}
%\SetWatermarkColor{red!30}

%\usepackage[printwatermark]{xwatermark}
%\newsavebox\mybox
%\savebox\mybox{\tikz[color=red,opacity=0.3]\node{\textsc{Entwurf}};}
%\newwatermark*[
%allpages,
%angle=45,
%scale=10,
%xpos=-4cm,
%ypos=4cm
%]{\usebox\mybox}
\pagestyle{fancy}
\fancyhead[L]{\includegraphics[width=2cm]{../images/logo_scaled.pdf}}
\fancyhead[R]{\textsc{Aufgabenserie 3}}



\begin{document}
%	\vspace*{-1cm}
%	\parbox{4cm}{\vspace{-0.2cm}\includegraphics[width=5cm]{../images/rolfal.jpg}}
%	\parbox{10.6cm}{\setstretch{2.0} \centering{ \huge \textsf{Aufgabenserie 3 
%			}}\\
%			Abgabe: 17. August \\ \vspace*{-.5cm} }
%		\vspace{0.5cm}
	\vspace*{-1cm}
	\parbox{4cm}{\vspace{-0.2cm}\includegraphics[width=5cm]{../images/logo_scaled.pdf}}
	\parbox{10.6cm}{\setstretch{2.0} \centering{ \huge \textsf{Aufgabenserie 3
			}}\\
			Abgabe: 17. August \\ \vspace*{-.5cm} }
		\vspace{0.5cm}

		
		\thispagestyle{empty}
\begin{framed}
	\noindent
	\scriptsize
	Diesmal gibt es vier Aufgaben, dafür auch zwei statt nur einem Monat. \textsc{ROLF wünscht viel Spaß in den Sommerferien.}\\
	Die Aufgaben sollten bis zum \textbf{17. August} bearbeitet werden. Die Lösungen schickt ihr an \href{mailto:physikrolf@gmail.com}{physikrolf@gmail.com}.
	Jede Aufgabe hat eine bestimmte Anzahl an erreichbaren Punkten. Wie viele das sind, müsst ihr raten. Versucht, die Lösungen so genau wie möglich aufzuschreiben. Für besonders schnelle/gute/witzige Lösungen kann es Bonuspunkte geben.\\ Die aktuellen Aufgaben sowie alle alten Aufgabenserien mit Lösungen findet ihr auch auf \url{pankratius.github.io/rolf}. \\\textit{Zu jeder Aufgabe gibt es jetzt Tipps. Die sollten beim Lösen der Aufgaben helfen.\\ Sollte das so sein macht bitte in euren Lösungen kenntlich, dass bestimmte Schritte von den Tipps und nicht von euch kommen. Darauf gibt es keinen Abzug, es ist nur für uns gut zu wissen.}
\end{framed}

\noindent

\begin{Exercise}[label = grassh, title = fauler Grashüpfer, origin = P.Gnädig, difficulty = 5]
	Ein fauler Grashüpfer möchte über einen Baumstumpf mit Radius $r = 20~\mathrm{cm}$ springen.\\
	Wie groß ist die dafür mindestens benötigte Geschwindigkeit, wenn der Luftwiderstand vernachlässigt werden kann?
\end{Exercise}
\begin{Answer}
	Bei minimaler Anfangsgeschwindigkeit sollte die Grashüpferbahn den Baumstamm in genau zwei Punkten symmetrisch zur Achse um den Baumstumpf berühren, und damit ihren Scheitelpunkt direkt auf dieser Achse haben, wie in Abbildung \ref{fig:grasshs} gezeigt ist. Dass die Bahn symmetrisch zur Symmetrieachse des Baums ist, liegt daran, dass sie reversibel sein sollte. Das bedeutet, dass es keinen Unterschied machen sollte, ob der Grashüpfer links oder rechts vom Baum losspringt, weshalb die Bahn symmetrisch sein muss. 
	%	\begin{figure}[h]
	%		\centering
	%		\begin{tikzpicture}[line cap=round,line join=round,>=triangle 45,x=1.0cm,y=1.0cm]
\clip(-3.573888715158432,-0.6284321850818549) rectangle (3.831338230521266,5.4106580654465875);
\draw [shift={(-3.,0.)}] (0,0) -- (0.:0.5107053065985999) arc (0.:72.38196924975453:0.5107053065985999) -- cycle;
\draw [shift={(-1.6166086789302696,3.1775297784800722)}] (0,0) -- (0.:0.5107053065985999) arc (0.:53.13793429316053:0.5107053065985999) -- cycle;
\draw [shift={(0.,2.)}] (0,0) -- (90.:0.5107053065985999) arc (90.:143.930590100419:0.5107053065985999) -- cycle;
\draw(0.,2.) circle (2.cm);
\draw [rotate around={90.:(0.,-25.3395997404539)}] (0.,-25.3395997404539) ellipse (29.719599740453898cm and 5.741326023857797cm);
\draw [->] (-1.6166086789302696,3.1775297784800722) -- (-0.76,4.32);
\draw (-1.6166086789302696,3.1775297784800722)-- (1.61660867893027,3.1775297784800722);
\draw (-1.6166086789302696,3.1775297784800722)-- (0.,2.);
\draw [->] (-3.,0.) -- (-2.6035486326210924,1.2484098166679993);
\begin{scriptsize}
\draw [fill=black] (-3.,0.) circle (2pt);
\draw [fill=black] (3.,0.) circle (2pt);
\draw [fill=black] (0.,2.) circle (2pt);
\draw [fill=black] (0.,0.) circle (2pt);
\draw [fill=black] (-1.6166086789302696,3.1775297784800722) circle (2pt);
\draw [fill=black] (1.61660867893027,3.1775297784800722) circle (2pt);
\draw [fill=black] (0.,4.38) circle (2pt);
%\draw [fill=black] (-0.76,4.32) circle (2.5pt);
\draw [fill=black] (0.,3.1775297784800722) circle (2pt);
%\draw [fill=black] (-2.6035486326210924,1.2484098166679993) circle (2.5pt);
\end{scriptsize}
\draw[thick] (-3.5,0)--(3.5,0);
\draw[thin,dashed](0,0)--(0,5);
\fill[pattern = north east lines] (-3.5,0)rectangle(3.5,-0.5);
\node at (-3.2,0.2){$A$};
\node at (3.2,0.2){$A^{\ast}$};
\node at (-2,3.1775297784800722){$B$};
\node at (2,3.1775297784800722){$B^{\ast}$};
\node at (0.5,2){$O$};
\node at (-1.25,2.5){$r$};
\node at (-2.25,0.2){$\alpha$};
\node at (-0.5,2.7) {$\theta$};
\node at (-0.9,3.4){$\theta$};
\node at (-1.8,3.8){$v_1$};
\node at (0.4,3.4){$C$};
\end{tikzpicture}
	%		\caption{Skizze der Grashüpferbahn}
	%		\label{fig:grasshs}		
	%	\end{figure}
	Es ist jetzt am einfachsten, wenn man zuerst die Bewegung des Grashüpfers von $B$ nach $B^{\ast}$ betrachtet. Die Geschwindigkeit des Grashüpfers in $B$ nennen wir $v_1$, und den Winkel, den die Grashüpfergeschwindigkeit mit dem Boden macht $\theta$.\\
	Dann wissen wir, dass die Wurfweite $s_w$ dieses schrägen Wurfes gegeben ist durch
	\begin{equation}\label{grassh:weite}
		s_w = \frac{2 v_1^2 \sin \theta \cos \theta}{g}.
	\end{equation}
	Geometrisch interpretiert muss diese Wurfweite $s_w$ genau die Strecke zwischen $B$ und $B^{\ast}$ sein. Wir versuchen jetzt, $\overline{BB^{\ast}}$ durch bekannte Parameter auszudrücken.\\
	Dazu macht es Sinn, wenn wir uns den Winkel $\angle COB$ anschauen. Wir wissen nämlich, dass die Geschwindigkeit (genauer: der Geschwindigkeits\textit{vektor}) tangiential zur Bahn sein muss, der Geschwindigkeitsvektor steht also senkrecht auf der Strecke $\overline{OB}$. Gleichzeitig ist das Dreieck $\Delta BOC$ ein rechtwinkliges, mit dem rechten Winkel $\angle BCO$. Damit ist also $\angle BOC = 180^\circ-90^\circ-\angle CBO = 90^\circ- \left(90^\circ- \theta\right) = \theta$.\\
	Jetzt können wir noch eine zweite Bedingung für $s_w$ aufstellen, allerdings in Abhängigkeit von $\theta$ und $r$. Das ist gut, weil wir dann einfach \eqref{grassh:weite} nehmen können, um $v_1$ auszurechen. Die Bedingung ist einfach, dass $s_w$ tatsächlich der Strecke $\overline{BB^\ast}$ entspricht. Es ist $\overline{BB^\ast} = 2 \overline{BC}$.\\
	Die Strecke $\overline{BC}$ können wir trigonometrisch durch $\theta$ im Dreieck $\Delta BOC$ ausrechnen. Dieses hat nämlich die Hypothenuse $r$, und der Winkel $\theta$ hat die Gegenkathete $\overline{BC}$. Also ist $\overline{BC} = r\cdot \sin \theta$. Damit haben wir jetzt die zweite Bedingung für $s_w$,
	\begin{equation}\label{grassh:weite2}
		s_w = \overline{BB^\ast} = 2 \overline{BC} = 2 r\cdot \sin \theta.
	\end{equation}
	Setzten wir \eqref{grassh:weite} und \eqref{grassh:weite2} gleich, kommen wir auf
	\begin{equation}\label{grassh:v2}
		\frac{2v_1^2\sin \theta \cos \theta}{g} = 2 r\cdot \sin \theta \Rightarrow v_1 = \sqrt{\frac{rg}{\cos\theta}}.
	\end{equation}
	Wir können jetzt die Energieerhaltung benutzten, um daraus die tatsächliche Anfangsgeschwindigkeit $v_0$ auszurechnen.\\
	Dazu brauchen wir aber noch die Höhe $h_B$ von Punkt $B$ überhalb des Bodens. Auch die können wir wieder durch $\theta$ ausrechnen. Die Höhe setzt sich nämlich zusammen aus der Höhe des Kreismittelpunkts $O$ über dem Boden (also dem Radius $r$) und der Strecke $\overline{OC}$. Die ist aber gerade die Ankathete zu $\theta$ im Dreieck $\Delta BOC$, sodass wir schreiben können $h_b = r + r\cos \theta = r\cdot\left(1+\cos \theta \right).$\\
	Die Energieerhaltung setzt nun die kinetische Energie am Punkt $A$ mit der Summe aus potentieller und kinetischer Energie in $B$ gleich, also
	\begin{equation}\label{grassh:coe}
		m\frac{v_0^2}{2} = m \frac{v_1^2}{2} + m g h_B\Rightarrow v_0^2 = \underbrace{\frac{rg}{\cos\theta}}_{\mathrm{aus}~\eqref{grassh:v2}} + 2g \underbrace{r\left(1+\cos \theta\right)}_{=h_B} =2 rg\cdot\left(1+\frac{1}{2\cos \theta}+ \cos \theta\right),
	\end{equation}
	wobei $m$ die Masse des Grashüpfers ist, die aber keine Rolle spielen sollte. \\
	Um den niedrigst möglichen  Wert für $v_0$ zu finden, reicht es, denn niedrigst möglichen Wert für $v_0^2$ zu finden. Mit \eqref{grassh:coe} geht das jetzt auf zwei unterschiedliche Wege. Für den einen kann man einfach ableiten und dann die Ableitung null setzten:
	\begin{equation}
		\frac{dv_0^2}{d\theta} \overset{!}{=} 0 \Rightarrow \frac{d}{d\theta}\left(\frac{1}{\cos \theta}+2\cos \theta\right)\overset{!}{=} 0\Rightarrow -2\sin \theta + \frac{1}{\cos^2 \theta} \sin \theta = 0 \Rightarrow \cos^2 \theta = \frac{1}{2} \Rightarrow \theta = 45^\circ.
	\end{equation}
	Hierbei wurde im ersten Schritt ausgenutzt, dass der erste Term in der Summe, aus der sich $v_0$ ergibt, nicht von $\theta$ abhängt, und somit für die Ableitung nicht relevant ist. Gleichzeitig kann der Faktor $rg$ gekürzt werden, weil $\frac{0}{rg} = 0$ gilt. Für die Ableitung von $\frac{1}{\cos \theta}$ kann man am einfachsten die Kettenregel nehmen. Das einfache Ergebnis für $\theta$ am Ende liegt daran, dass wir wissen, das $\theta>0$ und $\theta<90^\circ$ sein muss.\\
	Wenn man nicht ableiten kann, ist das aber nicht schlimm. Hier hilft die Ungleichung vom arithmetischen und geometrischen Mittel
	\begin{equation}\label{grassh:arge}
	\frac{x_1+x_2+...+x_n}{n}\geq \sqrt[n]{x_1\cdot x_2 \cdot ... \cdot x_n},
	\end{equation} 
	wobei man den linken Term arithmetisches Mittel (oft auch nur Durchschnitt) nennt, und den rechten geometrisches Mittel.
	Wir können jetzt den interessanten Term, also $2 \cos \theta + \frac{1}{\cos \theta}$ schreiben als 
	\begin{equation}\label{grassh:arge1}
		\frac{1}{2}\left(\frac{2}{\cos \theta} + 4 \cos \theta\right).
	\end{equation}
	In \eqref{grassh:arge1} haben wir den Term aber nur als arithmetisches Mittel zweier Terme ausgedrückt. Das muss aber immer mindestens genauso groß sein, wie das geometrische Mittel der beiden Teile, also
	\begin{equation}\label{grassh:arge2}
		\frac{1}{2}\left(\frac{2}{\cos \theta} + 4 \cos \theta\right) \geq \sqrt[2]{\frac{2}{\cos \theta}\cdot 4 \cos \theta} = \frac{\sqrt{2}}{2}.
	\end{equation}
	Denn kleinstmöglichen Wert für $v_0$ erhalten wir dann, wenn hier die Gleichheit gilt, weil ja den kleinstmöglichen Wert für die Summe aus \eqref{grassh:coe} darstellt. Stellen wir die Gleichung \eqref{grassh:arge2} nach $\theta$ um, so erhalten wir übrigens wieder $\theta = 45^\circ$.\\
	Setzt man das jetzt in \eqref{grassh:coe} ein, kommt man am Ende auf
	\begin{equation*}
		\boxed{
		v_0 = \sqrt{2gr\left(1+\sqrt{2}\right)}\approx 2.2~\mathrm{m\cdot s^{-1}}.}
	\end{equation*}
\end{Answer}

\textbf{Tipp 1:}
\textit{Bei minimaler Anfangsgeschwindigkeit ist die Bahn spiegelsymmetrisch zur Mittelpunktsachse des Baumstamms, und berührt ihn an insgesamt zwei Punkten (Warum?). Welche Kriterien lassen sich daraus ableiten?}
\\
\textbf{Tipp 2:}
\textit{Mit der Ungleichung vom arithmetischen und geometrischen Mittel kann man minimal oder maximal mögliche Werte finden:}
\begin{equation*}
	\frac{x_1+x_2+...+x_n}{n} \geq \sqrt[n]{x_1\cdot x_2 \cdot ... \cdot x_n}.
\end{equation*}

\begin{Exercise}[label = infnet, title = Widerstandsnetzwerk, origin = P.Gnädig, difficulty = 4]
	Alle Kanten in dem unendlich großen Widerstandsnetz (siehe Abbildung) haben den Widerstand $R$. Wie groß ist der Widerstand zwischen $A$ und $B$?
\end{Exercise}

\begin{figure}[h]
	\centering
	\begin{tikzpicture}
		\draw
		(-2,1) to[/tikz/circuitikz/bipoles/length=20pt, R] (0,1)
		(0,1) to[/tikz/circuitikz/bipoles/length=20pt, R](1,1)
		(1,1) to[/tikz/circuitikz/bipoles/length=20pt, R](3,1)
		(-1.75,1)to[/tikz/circuitikz/bipoles/length=20pt, R](-1.75,0)
		(-1.75,0)to[/tikz/circuitikz/bipoles/length=20pt, R](-1.75,-1)
				(-2,0) to[/tikz/circuitikz/bipoles/length=20pt, R] (0,0)
				(0,0) to[/tikz/circuitikz/bipoles/length=20pt, R](1,0)
				(1,0) to[/tikz/circuitikz/bipoles/length=20pt, R](3,0)
						(-2,-1) to[/tikz/circuitikz/bipoles/length=20pt, R] (0,-1)
						(0,-1) to[/tikz/circuitikz/bipoles/length=20pt, R](1,-1)
						(1,-1) to[/tikz/circuitikz/bipoles/length=20pt, R](3,-1)
				(-.25,1)to[/tikz/circuitikz/bipoles/length=20pt, R](-.25,0)
				(-.25,0)to[/tikz/circuitikz/bipoles/length=20pt, R](-.25,-1)
					(1.25,1)to[/tikz/circuitikz/bipoles/length=20pt, R](1.25,0)
					(1.25,0)to[/tikz/circuitikz/bipoles/length=20pt, R](1.25,-1)
					(2.75,1)to[/tikz/circuitikz/bipoles/length=20pt, R](2.75,0)
					(2.75,0)to[/tikz/circuitikz/bipoles/length=20pt, R](2.75,-1)
			
					;
		\draw (2.75,1) -- (2.75,1.2);
		\draw (-1.75,1) -- (-1.75,1.2);
		\draw (2.75,-1) -- (2.75,-1.2);
		\draw (-1.75,-1)--(-1.75,-1.2);
		\filldraw[black] (-.25,0) circle (1.5pt);
		\filldraw[black] (1.25,0) circle (1.5pt);
		\node at (0,-.2) {$A$};
		\node at (1.5,-.2) {$B$};
		\node at (0.5,0.25){$R$};
		\node at (-2.2,1) {$...$};
			\node at (-2.2,0) {$...$};
				\node at (-2.2,-1) {$...$};
		\node at (3.2,1) {$...$};
			\node at (3.2,0) {$...$};
			\node at (3.2,-1) {$...$};
		\node at (-1.75,-1.5) {$...$};
		\node at (-0.25,-1.5) {$...$};
		\node at (1.25,-1.5) {$...$};
		\node at (2.75,-1.5) {$...$};
				\node at (-1.75,1.5) {$...$};
				\node at (-0.25,1.5) {$...$};
				\node at (1.25,1.5) {$...$};
				\node at (2.75,1.5) {$...$};
		
		

	\end{tikzpicture}
\end{figure}
\textbf{Tipp 1:}
\textit{Mit dem Superpositionsprinzip aus Serie 2 kann man auch diese Aufgabe schön einfach lösen.}
\begin{minipage}[b]{0.6\textwidth}
\begin{Exercise}[label = col1, title = Stoßaufnahme, origin = Aaron Wild, difficulty = 5]
	Die nebenstehende Abbildung zeigt die Positionen von zwei Körpern zum Zeitpunkt $t_0$ und $t_1$. Die Körper sind sehr klein, und bewegen sich reibungsfrei auf einem Tisch.\\
	Der rote Körper hat eine Masse, die dreimal so groß ist wie die des blauen.\\
	Bestimme in der Abbildung (in größerer Form auch auf der nächsten Seite) die Bewegungsrichtungen der beiden Körper, nachdem sie zusammengestoßen sind.
\end{Exercise}
\end{minipage}
\begin{minipage}[b]{0.4\textwidth}
	\flushright
	\begin{tikzpicture}
	\filldraw[color =black, fill = blue]  (-2,2.31) circle (2pt);
	\filldraw[color = black, fill = blue] (0,1.155) circle (2pt);
	\filldraw[color = black, fill = red] (0,0) circle (2pt);
	\filldraw[color = black, fill = red] (1,0) circle (2pt);
	\node at (-2,2.7) {$t_0$};
	\node at (0,1.5) {$t_1$};
	\node at (0,-0.5) {$t_0$};
	\node at (1,-0.5){$t_1$};
\end{tikzpicture}
\end{minipage}
\textbf{Tipp 1:}
\textit{Offensichtlich bleibt der Gesamtimpuls und die Gesamtenergie erhalten. Gleichzeitig ändert sich aber für jeden Körper auch eine Impulskomponente nicht (Welche? Warum?). Damit  lässt sich der Rest einfacher konstruieren. }
\newpage
\begin{minipage}[b]{0.6\textwidth}
\begin{Exercise}[title = Variabler Brechungsindex, label = varn1, origin = IPhO 1974, difficulty = 4]
	Gegeben ist eine planparalle Platte. Ihre Brechzahl ändert sich nach der Gleichung
	\begin{equation}\label{varn1:nx}
		n\left(x\right) = \frac{n_a}{1-\frac{x}{q}},
	\end{equation}
	wobei $n_a = 1.2$ die Brechzahl im Punkt $A$, und $q = 0.13~\mathrm{cm}$ eine Konstante ist. \\
	Im Punkt $A$ ($x_a = 0$) fällt senkrecht zur Platte ein Lichstrahl ein. Dieser verlässt die Platte im Punkt $B$ unter einem Winkel von $\alpha = 30^{\circ}$ zur ursprünglichen Richtung. Bestimme den Brechungsindex des Materials im Punkt $B$ und die Dicke der Platte $d$.
\end{Exercise}
\end{minipage}
\begin{minipage}[b]{0.4\textwidth}
\centering
\begin{tikzpicture}
\clip (-1.5,-2.75) rectangle (4,4);
\draw[thick] (-0.5,1) -- (2.75,1);
\draw[thick] (-0.5,-1) -- (2.75,-1);
\draw[decorate,decoration={snake, amplitude = 0.7mm, segment length = 5.5mm}] (-0.5,1) -- (-0.5,-1);
\draw[decorate,decoration={snake, amplitude = 0.7mm, segment length = 5.5mm}] (2.5,1) -- (2.5,-1);
\draw[very thick,->] (0,-1) -- (0,2);
\draw[very thick,->] (0,-1) -- (3,-1);
\node at (-0.25,1.9) {$y$};
\node at (3.2,-1) {$x$};
\draw[<->] (-1,1) -- (-1,-1) node[midway, fill = white] {$d$};
	\begin{scope}[ every node/.style={sloped,allow upside down}
	] 
			\draw(0,-1.75) -- node {\midarrow} (0,-1);
	\end{scope}
\node at (-0.2,-1.2) {$A$};
\node at (1.5,0.7) {$B$};
\tkzDefPoint(1.5,1){B}
\tkzDefPoint(0,-1){A}
\tkzDefPoint(1.5,1.86){b1}
\tkzDefPoint(2,1.866025){b2}
\tkzMarkAngle[scale = 0.75](b2,B,b1)
\tkzLabelAngle(b2,B,b1){$\alpha$}
\tkzDrawSegment(B,b1)
\tkzDrawSegment(B,b2)
\tkzDefPoint(0.9,0.4){C}
\tkzCircumCenter(A,C,B)\tkzGetPoint{O}
\tkzDrawArc(O,B)(A)
\end{tikzpicture}
\end{minipage}
\begin{Answer}[ref = varn1]
	Wir rechnen zuerst den Brechungsindex des Materials im Punkt $B$ aus. Dazu hilft uns das Brechungsgesetz. Allgemein gilt für einen Strahl, der aus einem Material mit Brechungsindex $n_1$ unter einem Winkel $\alpha_1$ zum Lot auf einen Übergang zu einem Material mit Brechungsindex $n_2$ triftt
	\begin{equation*}
		n_1 \sin \alpha_1 = n_2 \sin \alpha_2,
	\end{equation*}
	wobei $\alpha_2$ der entsprechende Ausfallswinkel ist.\\
	Wir können uns das Material mit dem sich änderten Brechungsindex so vorstellen, als sei es aus sehr vielen kleinen Schichten mit Brechungsindex $n_i$ aufgebaut. Da jeweils das Brechungsgesetz gilt, muss also für zwei direkt hintereinader liegende Schichten gelten:
	\begin{equation*}
		n_{i-1} \sin \alpha_{i-1} = n_{i} \sin \alpha_i.
	\end{equation*}
	Das ist aber nichts anderes als die Aussage, dass $n \sin \alpha := c$ entlang des gesamten Materials konstant bleiben muss.\\
	Damit können wir jetzt den Brechungsindex im Punkt $B$ ausrechnen. Weil das Licht senkrecht eintrifft, ist
	\begin{equation}\label{varn1:cdef}
	c = n_a \sin 90^\circ = n_a.
	\end{equation}
	Gleichzeitig gilt für die Brechung am Punkt $B$ 
	\begin{equation}\label{varn1:bdef}
	\sin \alpha = n_B \sin \left(90^\circ-\beta\right) = n_B \cos \beta ,
	\end{equation}
	wobei $\beta_b$ der Einfallswinkel im Punkt $B$ ist. Mit dem trigonmetrische Pythagoras ($\sin^2 x + \cos^2 x = 1$) und $c = n_B \sin \beta_B$ kann man \eqref{varn1:bdef} umstellen zu
	\begin{multline}\label{varn1:nb}
		\sin \alpha = n_b \sqrt{1-\sin^2 \beta_B} =  \sqrt{n_b^2-n_b^2 \sin^2 \beta_B} = \sqrt{n_b^2-c^2} \overset{\eqref{varn1:cdef}}{=} \sqrt{n_b^2-n_a^2}\\ \boxed{\Rightarrow n_b = \sqrt{\sin^2 \alpha + n_a^2} = 1.3.}
	\end{multline}
	Wir können jetzt die Form des Lichstrahls ausrechnen, und damit $d$ berechnen. Das geht über mehrere Methoden.\\
	Bei der ersten denkt man ein bisschen nach. Dazu betrachtet man den Einfallswinkel $\beta\left(x\right)$ bei dem Übergang zwischen zwei gedachten Schichten an der Stelle $x$. Für diesen gilt mit \eqref{varn1:cdef} und \eqref{varn1:nx}
	\begin{equation}\label{varn1:circdef}
		\sin \beta\left(x\right) = \frac{c}{n\left(x\right)} \overset{\eqref{varn1:cdef}}{=} \frac{n_a}{n\left(x\right)} \overset{\eqref{varn1:nx}}{=} \frac{q-x}{q}.
	\end{equation}
	Wir betrachten nun den Kreis $k$ um den Punkt $O$, welcher die Koordinaten $\left(q,0\right)$ hat. Die Tangente an diesen Kreis an der Stelle $x$ schließt mit dem Kreis genau den Winkel $\sin \beta\left(x\right)$ ein, wie man im Dreieck $\Delta OCC'$ sehen kann. Weil der Strahl auch wenn er in das Material kommt tangential zu diesem Kreis verläuft, muss er innerhalb des Materials entlag dieser Kreisbahn verlaufen.\\
	Gleichzeitig können wir die $x$-Koordinate von $B$ einfach durch Umstellen von \eqref{varn1:nx} finden
	\begin{equation}\label{varn1:xb}
		n_b = \frac{n_a}{1-\frac{x}{q}}\Rightarrow x_B =q \left(1-\frac{n_a}{n_b}\right) = 1~\mathrm{cm}.
	\end{equation}
	Damit können wir jetzt im rechtwinkligen Dreieck $\Delta BB'O$ die Dicke $d$ ausrechnen
	\begin{equation}\label{varn1:dsol}
	\boxed{
		d = \sqrt{q^2-\left(q-x_b\right)^2} = 5~\mathrm{cm}.}
	\end{equation}
	Mann kann nachdenken auch durch Ableiten ersetzen. Mit $\tan\left(\arcsin x\right) = \frac{x}{\sqrt{1-x^2}}$ können wir aus \eqref{varn1:circdef} den Anstieg der Funktion $y\left(x\right)$ für den Lichtstrahl bestimmen
	\begin{equation*}
		\frac{\mathrm{d}y}{\mathrm{d}x} = \tan \beta = \frac{q-x}{\sqrt{q^2-\left(q-x\right)^2}}.
	\end{equation*}
	Das ist eine separierbare Differentialgleichung, die wir integrieren können
	\begin{equation}\label{varn1:inte}
		\int \mathrm{d}y = \int \frac{q-x}{\sqrt{q^2-\left(q-x\right)^2}}~\mathrm{d}x \Rightarrow y\left(x\right)= \sqrt{q^2-\left(q-x\right)^2}+c
	\end{equation}
	Aus $y\left(0\right)= 0$ folgt $c=0$. \eqref{varn1:inte} ist aber nur die Gleichung für einen Halbkreis in positiver $y$-Richtung um $\left(q,0\right)$. Damit kann $d$ genauso wie in \eqref{varn1:dsol} ausgerechnet werden.
	
\end{Answer}
\textbf{Tipp 1:}
\textit{Es gilt das Brechungsgesetz (Tafelwerk!) für den Übergang zwischen Schichten mit unterschiedlichem Brechungsindex. Lässt sich daraus vielleicht eine Erhaltungsgröße ableiten?}

\newpage
	\hspace{4cm}
	\vspace{-1cm}
\begin{turn}{270}

\begin{tikzpicture}[scale = 3]
	\filldraw[color =black, fill = blue]  (-2,12.31) circle (0.05);
	\filldraw[color = black, fill = blue] (0,11.155) circle (0.05);
	\filldraw[color = black, fill = red] (-0.1,10) circle (0.05);
	\filldraw[color = black, fill = red] (0.9,10) circle (0.05);

\end{tikzpicture}
\end{turn}



\end{document}
